\documentclass{oxmathproblems}

\printanswers

\course{Understanding Analysis, 2nd Edition, Stephen Abbott}
\sheettitle{Chapter 4 - Section 4.4 - Continuous Functions on Compact satisfies} 

\begin{document}
\begin{questions}

%------------------------- Problem 1 -------------------------
\miquestion 
\begin{enumerate}[label=(\alph*)]
  \item Show that $f(x)=x^3$ is continuous on all of \textbf{R}.
  \item Argue, using Theorem 4.4.5, that $f$ is not uniformly continuous on \textbf{R}.
  \item Show that $f$ is uniformly continuous on any bounded subset of \textbf{R}.
\end{enumerate}

\begin{solution}
  \begin{enumerate}[label=(\alph*)]
    \item $f(x)=x$ is continuous on all of \textbf{R}. Product of continuous functions is continuous, so $f(x)=x^3$ is continuous as well.
    \item As per the Theorem, we need a sequence $(x_{n})$, $(y_{n})$ such that $|x_{n}-y_{n}| \to 0$ but $|f(x_{n})-f(y_{n})| \geq \epsilon$ for a choice of $\epsilon$.
    We'll choose $\epsilon=1$. Intuitively, as $x$ and $y$ gets closer we place them farther and farther out where $x^3$ becomes extremely steep.

    We'll construct a sequence such that $|x_{n}-y_{n}| = \frac{1}{n}$ for all $n$. We have $|f(x)-f(y)|=|x^3-y^3|=|x-y||x^2+xy+y^2|$.
    Since $|x-y| = \frac{1}{n}$, to make the expression $\geq 1$ we need $|x^2+xy+y^2| \geq n$. Simply choose $x, y \geq n$ for each $n$,
    then $|x^2+xy+y^2| > n^2 > n$, as required. One possible sequence is $(x_{n})=n$ and $(y_{n})=n+\frac{1}{n}$.

    \item Since the subset is bounded, there exists an $M > 0$ such $|x| \leq M$ for all $x$. The idea is that the delta for $M$ works for all other $|x| < M$ as well, 
    because $x^3$ gets steeper farther out from zero and requires smaller and smaller $\delta$.

    More formally, do the usual $\epsilon-\delta$ proof to get $|x-y||x^2+xy+y^2|$. We have $|x|, |y| \leq M$ so $|x-y||x^2+xy+y^2| \leq 3M^2|x-y|$. Pick $\delta=\frac{\epsilon}{3M^2}$.
    For an $\epsilon > 0$, this single $\delta$ works for all $x$.
  \end{enumerate}
\end{solution}

%------------------------- Problem 2 -------------------------
\miquestion
\begin{enumerate}[label=(\alph*)]
  \item Is $f(x)=1/x$ uniformly continuous on $(0, 1)$?
  \item Is $g(x)=\sqrt{x^2+1}$ uniformly continuous on $(0, 1)$?
  \item Is $h(x)=x \sin(1/x)$ uniformly continuous on $(0, 1)$?
\end{enumerate}

\begin{solution}
  \begin{enumerate}[label=(\alph*)]
    \item No, as it gets infinitely steeper closer to zero. We'll prove it like problem 1 part (b) and construct $(x_{n})$, $(y_{n})$ such that
    $|x_{n}-y_{n}| \to 0$ but $|f(x_{n})-f(y_{n})| \geq 1$.

    Construct the sequences such that $|x_{n}-y_{n}|=\frac{1}{n}$ for all $n$. We have $|f(x)-f(y)|=|x-y|\frac{1}{|xy|}$. To make this greater than one, we need
    $\frac{1}{|xy|} \geq n \implies |xy| \leq \frac{1}{n}$. One such sequence is $(x_{n})=\frac{1}{n}$ and $y_{n}=\frac{2}{n}$, starting from $n = 3$ so that it
    satisfies the requirement that the domain is $(0, 1)$.

    \item Yes. Do the usual $\epsilon-\delta$ to get
    \begin{align*}
       |\sqrt{x^2+1}-\sqrt{y^2+1}| &= \frac{|x^2-y^2|}{\sqrt{x^2+1}+\sqrt{y^2+1}}\\
       &= |x-y| \frac{|x+y|}{\sqrt{x^2+1}+\sqrt{y^2+1}}
    \end{align*}

    We need to produce a $\delta$ so that the above is less than $\epsilon$. Since $x$ and $y$ are bounded by $(0, 1)$, the $\frac{|x+y|}{\sqrt{x^2+1}+\sqrt{y^2+1}}$ term
    is bounded by some constant $c$ that we don't need to bother computing; the important thing is that it's a constant. Set $\delta=\frac{\epsilon}{c}$, then
    \[|x-y| \frac{|x+y|}{\sqrt{x^2+1}+\sqrt{y^2+1}} < \frac{\epsilon}{c} \cdot c = \epsilon\]
    This delta works for all $y$ for a choice of $\epsilon$.

    Or if we want to compute the actual bound. We have $x, y < 1$ so $|x+y| < 2$. Also, $x, y > 0$ so $\sqrt{x^2+1}+\sqrt{y^2+1} > 2$. Then $\frac{|x+y|}{\sqrt{x^2+1}+\sqrt{y^2+1}} < 1$,
    so we can set $\delta=\epsilon$. Then
    \[|x-y| \frac{|x+y|}{\sqrt{x^2+1}+\sqrt{y^2+1}} < |x-y| < \epsilon\]

    \item Looking at the graph, the answer could possibly be yes. It has a similar problem as $\sin(1/x)$ in that it oscillates wildly closer to zero, but otoh the amplitude gets compressed, so
    each oscillation covers smaller range of values.

    Let's try to prove it's possible and see where it goes. Consider $|x \sin(1/x) - y \sin(1/y)|$. By triangle inequality and the fact that $\sin(1/x) \in [-1, 1]$ for $x \in (0, 1)$, we have
    \[|x \sin(1/x) - y \sin(1/y)| \leq |x \sin(1/x)| + |y \sin(1/y)| \leq x + y\]
    This bound is only useful if we limit $x, y$ to be small enough. Say, one of $x, y < \epsilon/4$ and $|x-y| < \epsilon/4$, which ensures $x+y < \epsilon/4 + \epsilon/2 < \epsilon$.

    What if $x, y \in [\epsilon/4, 1)$? Well, $x \sin(1/x)$ is continuous on $[\epsilon/4, 1]$ (note the closed interval). By Theorem 4.4.7, it's uniformly continuous on $[\epsilon/4, 1]$ and since
    $[\epsilon/4, 1) \subset [\epsilon/4, 1]$, it's uniformly continuous on $[\epsilon/4, 1)$ as well. Hence there's a single $\delta_{1}$ that works for this region.

    All in all, choose $\delta = \min(\delta_{1}, \epsilon/4)$.
    
    --- Much more elegant proof by chatgpt ---\\
    Extend $h(x) = x \sin(1/x)$ to 0 by defining $h(0)=0$. Now $h(x)$ is continuous on [0, 1], hence uniformly continuous on [0, 1], hence uniformly continuous on (0, 1). This is in fact using the result
    of Exercise 4.4.13 (Continuous Extension Theorem).
  \end{enumerate}
\end{solution}

\end{questions}
\end{document}
